\section{Introduction}

Large Language Model (LLM) agents can use natural-language instructions, external tools, and execution environments to complete complex multi-step tasks. As agent capabilities advance, method descriptions in research papers are increasingly treated as reusable sources of procedural knowledge. Existing studies use agents to recreate code and experiments described in papers \cite{starace2025paperbench,zhao2025autoreproduce}, while paper-to-agent systems further transform papers and associated software into interactive agent artifacts \cite{miao2025paper2agent}. Through such upstream conversion processes, paper methods can be organized into executable agent skill specifications.

Existing research follows two related directions. The first extracts method content from papers or other knowledge sources and produces executable agent artifacts or reusable skills \cite{miao2025paper2agent,shen2026skillfoundry,pan2026anything2skill}. The second reduces context length and execution cost through prompt compression or agent-skill compression \cite{jiang2023llmlingua,jiang2024longllmlingua,pan2024llmlingua2,gao2026skillreducer}. For example, LLMLingua removes tokens estimated to carry limited information, LongLLMLingua introduces question-aware compression for long contexts, and LLMLingua-2 formulates compression as token-level classification. However, removing tokens or sentences can disrupt semantic coherence and discard necessary information at high compression ratios \cite{pan2024llmlingua2,liskavets2025contextcompression}.

A shorter skill specification does not necessarily retain all information required for method reproduction. Existing methods commonly evaluate simplification using compression ratio, execution cost, or downstream task performance, but rarely jointly record which source passages were removed, whether the remaining steps preserve their prerequisites, and which experimental results justify accepting or rejecting a candidate. Consequently, success on a small set of examples does not distinguish an insufficient full specification, a simplified candidate missing critical information, and an invalid experimental run. We therefore study candidate-level post-compression validation: given a full skill specification, a simplified candidate specification, and an executable evaluation package, determine whether the candidate can substitute for the full specification for a specified task under fixed experimental conditions.

We present SkillAudit, a traceable post-compression validation protocol. SkillAudit takes a source-to-step mapping confirmed by the evaluation-package builder, binds procedural steps to source passages and prerequisites, and checks whether a simplified candidate preserves the dependency closure required for task execution. Before held-out evaluation, the protocol freezes the candidate, task, workspace, model, tools, run budget, tests, scorer, and prespecified acceptance criteria. It then executes three conditions: a no-specification baseline, the full specification, and the simplified candidate specification. Finally, SkillAudit combines held-out results with integrity checks, separates candidate rejection from baseline-, full-, and evidence-inconclusive states, and emits an auditable record containing source provenance, experimental configuration, results, and the decision rationale.

We evaluate matched evidence with GPT-5.6 Sol as the homogeneous primary executor. Three fresh schedule-matched repetitions yield \PrimaryRows{} terminal rows from \PrimaryPapers{} papers and \PrimaryTasks{} paper-task units spanning NLP, software engineering, data analysis, and agent/tool use. Relative to the no-specification baseline, $F$ and $S$ improve operational success by \FBSuccessEffect{} and \SBSuccessEffect{}, respectively, with positive paper-clustered intervals; their $S-F$ difference is \SFSuccessEffect{} and its interval crosses zero. Both baseline contrasts remain positive in every repetition and domain, whereas the local $S-F$ ordering varies. A separate known-truth suite recovers all \ControlDeterministicCells{} deterministic labels and distinguishes identity or redundancy-preserving changes from deletion of a registered exact-output contract. Because the six-state policy was not bound before the primary executions, its pooled replay remains retrospective rather than confirmatory admission evidence. Together, the experiments support $B/F/S$ comparison and integrity-aware controls as task-local audit evidence; they do not establish candidate equivalence or universal admissibility.

Our contributions are threefold:
\begin{itemize}
    \item We formulate candidate-level post-compression validation and distinguish failures caused by an insufficient full specification, missing information in a simplified candidate, baseline solvability, and invalid experimental runs.
    \item We introduce SkillAudit, which integrates source mapping, dependency-closure checking, controlled three-condition execution, a frozen evaluation boundary, and auditable decision records into a unified validation protocol.
    \item We provide matched evidence through three fresh GPT-5.6 Sol repetitions over \PrimaryTasks{} paper-task units and \PrimaryDomains{} domains, paper-clustered effects, known-truth controls, component ablation, and bounded robustness analyses, while exposing where the current registry remains inconclusive.
\end{itemize}
